% ============================================================
% EJERCICIOS
% CARGA ELÉCTRICA Y ELECTRIZACIÓN
% ============================================================

\section{Ejercicios}


% ============================================================
% NIVEL BÁSICO
% ============================================================

\subsection{Nivel básico}

Los ejercicios de este nivel buscan reconocer los conceptos fundamentales,
interpretar el signo de la carga, aplicar la cuantización y distinguir los
principales procesos de electrización.


% ------------------------------------------------------------
% BÁSICO 1
% ------------------------------------------------------------

\begin{ejercicio}

Un cuerpo inicialmente neutro pierde \(3.0\times10^{12}\) electrones.

Determine:

\begin{enumerate}
    \item el signo de la carga final;
    \item la magnitud de la carga adquirida.
\end{enumerate}

\end{ejercicio}

\begin{respuesta}

El cuerpo queda cargado positivamente y

\[
\boxed{
Q\approx+\SI{4.81e-7}{\coulomb}
}
\]

o

\[
\boxed{
Q\approx+\SI{0.481}{\micro\coulomb}
}.
\]

\end{respuesta}

\begin{explicacion}

Como el cuerpo \textbf{pierde electrones}, pierde carga negativa.
Por tanto, queda con un déficit de electrones y su carga final debe ser
positiva:

\[
Q>0.
\]

La carga correspondiente a \(N\) electrones tiene magnitud

\[
|Q|=Ne.
\]

Como en este caso los electrones fueron perdidos,

\[
Q=+Ne.
\]

Sustituyendo,

\[
Q=
\left(3.0\times10^{12}\right)
\left(\SI{1.602176634e-19}{\coulomb}\right).
\]

Entonces,

\[
Q\approx\SI{4.81e-7}{\coulomb}.
\]

Por tanto,

\[
\boxed{
Q\approx+\SI{4.81e-7}{\coulomb}
}
\]

o

\[
\boxed{
Q\approx+\SI{0.481}{\micro\coulomb}
}.
\]

El signo positivo representa físicamente un déficit de electrones.

\end{explicacion}


% ------------------------------------------------------------
% BÁSICO 2
% ------------------------------------------------------------

\begin{ejercicio}

Un cuerpo inicialmente neutro gana \(8.0\times10^{11}\) electrones.

Determine su carga final.

\end{ejercicio}

\begin{respuesta}

\[
\boxed{
Q\approx-\SI{1.28e-7}{\coulomb}
}
\]

o

\[
\boxed{
Q\approx-\SI{0.128}{\micro\coulomb}
}.
\]

\end{respuesta}

\begin{explicacion}

Al ganar electrones, el cuerpo adquiere un exceso de carga negativa.
Por tanto,

\[
Q<0.
\]

La carga adquirida es

\[
Q=-Ne.
\]

Sustituyendo,

\[
Q=
-\left(8.0\times10^{11}\right)
\left(\SI{1.602176634e-19}{\coulomb}\right).
\]

Entonces,

\[
Q\approx-\SI{1.28e-7}{\coulomb}.
\]

Expresando el resultado en microcoulombs,

\[
\boxed{
Q\approx-\SI{0.128}{\micro\coulomb}
}.
\]

El signo negativo indica que el cuerpo posee electrones en exceso.

\end{explicacion}


% ------------------------------------------------------------
% BÁSICO 3
% ------------------------------------------------------------

\begin{ejercicio}

Un cuerpo posee una carga

\[
Q=-\SI{3.204353268e-18}{\coulomb}.
\]

¿Cuántos electrones en exceso posee?

\end{ejercicio}

\begin{respuesta}

\[
\boxed{
N=20\text{ electrones}
}.
\]

\end{respuesta}

\begin{explicacion}

El número de electrones asociados con una determinada magnitud de carga se
calcula mediante

\[
N=\frac{\abs{Q}}{e}.
\]

Sustituyendo,

\[
N=
\frac{3.204353268\times10^{-18}}
{1.602176634\times10^{-19}}.
\]

Por tanto,

\[
N=20.
\]

Como la carga del cuerpo es negativa, esos electrones representan un exceso
respecto del estado neutro.

\[
\boxed{
N=20\text{ electrones en exceso}
}.
\]

El número \(N\) no se escribe con signo negativo porque representa una
cantidad de partículas.

\end{explicacion}


% ------------------------------------------------------------
% BÁSICO 4
% ------------------------------------------------------------

\begin{ejercicio}

Indique si cada cuerpo es positivo, negativo o neutro.

\begin{enumerate}
    \item Tiene el mismo número de protones y electrones.
    \item Tiene \(15\) electrones más que protones.
    \item Tiene \(8\) electrones menos que protones.
\end{enumerate}

\end{ejercicio}

\begin{respuesta}

\begin{enumerate}
    \item Neutro.
    \item Negativo.
    \item Positivo.
\end{enumerate}

\end{respuesta}

\begin{explicacion}

Un cuerpo es neutro cuando sus cargas positivas y negativas se compensan.

\begin{enumerate}

    \item Si posee igual número de protones y electrones,

    \[
    Q=0,
    \]

    por lo que es eléctricamente neutro.

    \item Si posee más electrones que protones, existe un exceso de carga
    negativa. Por tanto, el cuerpo es negativo.

    \item Si posee menos electrones que protones, existe un déficit de
    electrones. Por tanto, el cuerpo es positivo.

\end{enumerate}

En los procesos electrostáticos ordinarios de sólidos, un cuerpo positivo
no suele haber ganado protones; normalmente ha perdido electrones.

\end{explicacion}


% ------------------------------------------------------------
% BÁSICO 5
% ------------------------------------------------------------

\begin{ejercicio}

Dos cuerpos poseen cargas

\[
Q_1=+\SI{4}{\micro\coulomb},
\qquad
Q_2=-\SI{7}{\micro\coulomb}.
\]

Determine la carga total del sistema.

\end{ejercicio}

\begin{respuesta}

\[
\boxed{
Q_{\text{total}}=-\SI{3}{\micro\coulomb}
}.
\]

\end{respuesta}

\begin{explicacion}

La carga eléctrica es una magnitud escalar, por lo que las cargas se suman
algebraicamente:

\[
Q_{\text{total}}=Q_1+Q_2.
\]

Sustituyendo,

\[
Q_{\text{total}}
=
\SI{4}{\micro\coulomb}
-\SI{7}{\micro\coulomb}.
\]

Por tanto,

\[
\boxed{
Q_{\text{total}}=-\SI{3}{\micro\coulomb}
}.
\]

El signo negativo indica que, en conjunto, el sistema presenta un exceso
neto de carga negativa.

\end{explicacion}


% ------------------------------------------------------------
% BÁSICO 6
% ------------------------------------------------------------

\begin{ejercicio}

Indique qué tipo de interacción cualitativa se espera entre las siguientes
parejas de cargas:

\begin{enumerate}
    \item positiva y positiva;
    \item negativa y negativa;
    \item positiva y negativa.
\end{enumerate}

\end{ejercicio}

\begin{respuesta}

\begin{enumerate}
    \item Repulsión.
    \item Repulsión.
    \item Atracción.
\end{enumerate}

\end{respuesta}

\begin{explicacion}

Las cargas del mismo signo se repelen y las cargas de signos opuestos se
atraen.

Por tanto,

\[
(+,+)\Longrightarrow\text{repulsión},
\]

\[
(-,-)\Longrightarrow\text{repulsión},
\]

mientras que

\[
(+,-)\Longrightarrow\text{atracción}.
\]

En este ejercicio solamente se analiza el carácter atractivo o repulsivo de
la interacción; no se calcula todavía su magnitud.

\end{explicacion}


% ------------------------------------------------------------
% BÁSICO 7
% ------------------------------------------------------------

\begin{ejercicio}

Clasifique cada material, en condiciones ordinarias, como conductor o
aislante:

\begin{enumerate}
    \item cobre;
    \item vidrio;
    \item aluminio;
    \item caucho;
    \item plástico.
\end{enumerate}

\end{ejercicio}

\begin{respuesta}

Conductores:

\[
\boxed{\text{cobre y aluminio}}
\]

Aislantes:

\[
\boxed{\text{vidrio, caucho y plástico}}
\]

\end{respuesta}

\begin{explicacion}

En los conductores existen portadores de carga que pueden desplazarse con
relativa facilidad a través del material.

Los metales, como el cobre y el aluminio, presentan esta característica.

En materiales como el vidrio, el caucho y muchos plásticos, las cargas poseen
mucha menor libertad de movimiento macroscópico.

Por ello, en condiciones ordinarias, estos materiales se comportan como
aislantes.

La diferencia no está en que unos materiales tengan electrones y otros no,
sino en la movilidad de sus portadores de carga.

\end{explicacion}


% ------------------------------------------------------------
% BÁSICO 8
% ------------------------------------------------------------

\begin{ejercicio}

Complete correctamente:

\begin{enumerate}
    \item Si un cuerpo gana electrones, queda cargado
    \underline{\hspace{2cm}}.
    \item Si un cuerpo pierde electrones, queda cargado
    \underline{\hspace{2cm}}.
    \item Si sus cargas positivas y negativas se compensan, el cuerpo es
    \underline{\hspace{2cm}}.
\end{enumerate}

\end{ejercicio}

\begin{respuesta}

\begin{enumerate}
    \item negativamente;
    \item positivamente;
    \item neutro.
\end{enumerate}

\end{respuesta}

\begin{explicacion}

Los electrones poseen carga negativa.

Por ello,

\[
\text{ganar electrones}
\Longrightarrow
\text{carga negativa},
\]

mientras que

\[
\text{perder electrones}
\Longrightarrow
\text{carga positiva}.
\]

Un cuerpo es eléctricamente neutro cuando la suma algebraica de todas sus
cargas es cero:

\[
Q=0.
\]

\end{explicacion}


% ------------------------------------------------------------
% BÁSICO 9
% ------------------------------------------------------------

\begin{ejercicio}

Indique el método de electrización descrito en cada situación:

\begin{enumerate}
    \item Una barra plástica se frota con una tela.
    \item Una esfera neutra toca otra esfera cargada.
    \item Un conductor se carga utilizando un cuerpo cercano sin tocarlo y
    una conexión a tierra.
\end{enumerate}

\end{ejercicio}

\begin{respuesta}

\begin{enumerate}
    \item Frotamiento.
    \item Contacto.
    \item Inducción.
\end{enumerate}

\end{respuesta}

\begin{explicacion}

En el primer caso, dos materiales se ponen en contacto y se frotan, lo que
puede favorecer la transferencia de electrones. Se trata de electrización
por frotamiento.

En el segundo caso, un cuerpo cargado toca directamente al cuerpo inicialmente
neutro y se produce transferencia de carga. Se trata de electrización por
contacto.

En el tercer caso, el cuerpo cargado que produce la influencia no toca el
conductor. La redistribución inicial de cargas y la conexión a tierra permiten
que el conductor adquiera carga neta. Se trata de electrización por inducción.

\end{explicacion}


% ------------------------------------------------------------
% BÁSICO 10
% ------------------------------------------------------------

\begin{ejercicio}

Un conductor inicialmente neutro se polariza por la presencia de un cuerpo
cargado cercano.

¿Cuál es su carga neta mientras no exista transferencia de carga con el
exterior?

\end{ejercicio}

\begin{respuesta}

\[
\boxed{
Q_{\text{neto}}=0
}.
\]

\end{respuesta}

\begin{explicacion}

La polarización produce una redistribución de las cargas que ya se
encuentran en el conductor.

Una región puede presentar un exceso relativo de electrones y otra un déficit,
pero ninguna carga ha entrado o salido del cuerpo.

Por tanto,

\[
\boxed{
Q_{\text{neto}}=0
}.
\]

El conductor está polarizado, pero continúa siendo eléctricamente neutro.

\end{explicacion}


% ============================================================
% NIVEL INTERMEDIO
% ============================================================

\subsection{Nivel intermedio}

En este nivel se combinan conservación, cuantización, transferencia de
electrones y análisis de procesos de electrización.


% ------------------------------------------------------------
% INTERMEDIO 1
% ------------------------------------------------------------

\begin{ejercicio}

Dos cuerpos inicialmente neutros se frotan entre sí. Después del proceso,
el cuerpo \(A\) posee una carga

\[
Q_A=+\SI{2.4}{\micro\coulomb}.
\]

Si el sistema formado por ambos cuerpos puede considerarse aislado, determine
la carga final del cuerpo \(B\).

\end{ejercicio}

\begin{respuesta}

\[
\boxed{
Q_B=-\SI{2.4}{\micro\coulomb}
}.
\]

\end{respuesta}

\begin{explicacion}

Inicialmente ambos cuerpos eran neutros, por lo que la carga total era

\[
Q_{\text{total,inicial}}=0.
\]

Como el sistema se considera aislado, se cumple la conservación de la carga:

\[
Q_{\text{total,inicial}}
=
Q_{\text{total,final}}.
\]

Por tanto,

\[
Q_A+Q_B=0.
\]

Despejando,

\[
Q_B=-Q_A.
\]

Sustituyendo,

\[
\boxed{
Q_B=-\SI{2.4}{\micro\coulomb}
}.
\]

Las cargas tienen igual magnitud y signos opuestos porque toda la carga
transferida de un cuerpo apareció en el otro y no hubo intercambio con el
exterior.

\end{explicacion}


% ------------------------------------------------------------
% INTERMEDIO 2
% ------------------------------------------------------------

\begin{ejercicio}

Un cuerpo inicialmente neutro adquiere una carga

\[
Q=-\SI{4.0}{\nano\coulomb}.
\]

Determine aproximadamente cuántos electrones ganó.

\end{ejercicio}

\begin{respuesta}

\[
\boxed{
N\approx2.50\times10^{10}\text{ electrones}
}.
\]

\end{respuesta}

\begin{explicacion}

Primero expresamos la carga en coulombs:

\[
Q=-\SI{4.0e-9}{\coulomb}.
\]

El número de electrones transferidos se obtiene mediante

\[
N=\frac{\abs{Q}}{e}.
\]

Sustituyendo,

\[
N=
\frac{4.0\times10^{-9}}
{1.602176634\times10^{-19}}.
\]

Entonces,

\[
N\approx2.50\times10^{10}.
\]

Como la carga final es negativa, el cuerpo ganó electrones.

Por tanto,

\[
\boxed{
N\approx2.50\times10^{10}\text{ electrones ganados}
}.
\]

\end{explicacion}


% ------------------------------------------------------------
% INTERMEDIO 3
% ------------------------------------------------------------

\begin{ejercicio}

Tres cuerpos poseen inicialmente las cargas

\[
Q_A=+\SI{5}{\micro\coulomb},
\qquad
Q_B=-\SI{2}{\micro\coulomb},
\qquad
Q_C=-\SI{1}{\micro\coulomb}.
\]

Después de interactuar dentro de un sistema aislado se encuentra que

\[
Q_A'=+\SI{1}{\micro\coulomb},
\qquad
Q_B'=-\SI{4}{\micro\coulomb}.
\]

Determine \(Q_C'\).

\end{ejercicio}

\begin{respuesta}

\[
\boxed{
Q_C'=+\SI{5}{\micro\coulomb}
}.
\]

\end{respuesta}

\begin{explicacion}

Calculamos primero la carga total inicial:

\[
Q_{\text{inicial}}
=
Q_A+Q_B+Q_C.
\]

Sustituyendo,

\[
Q_{\text{inicial}}
=
(5-2-1)\,\si{\micro\coulomb}.
\]

Entonces,

\[
Q_{\text{inicial}}
=
+\SI{2}{\micro\coulomb}.
\]

Por conservación de la carga,

\[
Q_{\text{final}}
=
+\SI{2}{\micro\coulomb}.
\]

Por tanto,

\[
Q_A'+Q_B'+Q_C'
=
+\SI{2}{\micro\coulomb}.
\]

Sustituyendo,

\[
1-4+Q_C'=2.
\]

Despejando,

\[
Q_C'=5.
\]

Así,

\[
\boxed{
Q_C'=+\SI{5}{\micro\coulomb}
}.
\]

La conservación se aplica al sistema completo, independientemente de cómo se
haya distribuido internamente la carga.

\end{explicacion}


% ------------------------------------------------------------
% INTERMEDIO 4
% ------------------------------------------------------------

\begin{ejercicio}

Determine cuáles de las siguientes cargas son compatibles con la
cuantización de la carga:

\begin{enumerate}
    \item \(Q=5e\);
    \item \(Q=-12e\);
    \item \(Q=3.5e\);
    \item \(Q=0\).
\end{enumerate}

\end{ejercicio}

\begin{respuesta}

Son compatibles:

\[
\boxed{
5e,\qquad -12e,\qquad 0
}.
\]

La carga

\[
\boxed{
3.5e
}
\]

no es compatible con el modelo estudiado.

\end{respuesta}

\begin{explicacion}

La cuantización de la carga establece

\[
Q=ne,
\]

donde

\[
n\in\mathbb{Z}.
\]

Los valores

\[
5,\qquad -12,\qquad 0
\]

son números enteros.

En cambio,

\[
3.5\notin\mathbb{Z}.
\]

Por tanto,

\[
Q=3.5e
\]

no corresponde a una carga neta permitida dentro del modelo electrostático
ordinario utilizado en este tema.

\end{explicacion}


% ------------------------------------------------------------
% INTERMEDIO 5
% ------------------------------------------------------------

\begin{ejercicio}

Dos esferas conductoras idénticas y aisladas poseen inicialmente

\[
Q_1=+\SI{8}{\micro\coulomb},
\qquad
Q_2=-\SI{2}{\micro\coulomb}.
\]

Se ponen en contacto y luego se separan.

Determine la carga final de cada esfera.

\end{ejercicio}

\begin{respuesta}

\[
\boxed{
Q_1'=Q_2'=+\SI{3}{\micro\coulomb}
}.
\]

\end{respuesta}

\begin{explicacion}

La carga total inicial es

\[
Q_{\text{total}}
=
Q_1+Q_2.
\]

Sustituyendo,

\[
Q_{\text{total}}
=
\SI{8}{\micro\coulomb}
-\SI{2}{\micro\coulomb}.
\]

Por tanto,

\[
Q_{\text{total}}
=
+\SI{6}{\micro\coulomb}.
\]

La carga se conserva.

Como las esferas son conductoras, idénticas y se encuentran en condiciones
simétricas, después del contacto quedan con la misma carga:

\[
Q_1'=Q_2'.
\]

Entonces,

\[
Q_1'=Q_2'
=
\frac{Q_{\text{total}}}{2}.
\]

Sustituyendo,

\[
Q_1'=Q_2'
=
\frac{\SI{6}{\micro\coulomb}}{2}.
\]

Finalmente,

\[
\boxed{
Q_1'=Q_2'=+\SI{3}{\micro\coulomb}
}.
\]

La división por partes iguales no es una regla general para cualquier par de
conductores; aquí es válida porque las dos esferas son idénticas.

\end{explicacion}


% ------------------------------------------------------------
% INTERMEDIO 6
% ------------------------------------------------------------

\begin{ejercicio}

Una esfera conductora neutra se encuentra cerca de una barra negativamente
cargada, sin que exista contacto.

Describa qué ocurre con los electrones de la esfera y determine su carga
neta.

\end{ejercicio}

\begin{respuesta}

Los electrones se desplazan hacia la región más alejada de la barra.

La esfera se polariza, pero

\[
\boxed{
Q_{\text{esfera}}=0
}.
\]

\end{respuesta}

\begin{explicacion}

La barra negativa repele a los electrones móviles del conductor.

Como consecuencia, estos se desplazan hacia la zona de la esfera más alejada
de la barra.

La región cercana queda con déficit relativo de electrones y la región
alejada con exceso.

Sin embargo, ningún electrón ha salido de la esfera ni ha ingresado desde el
exterior.

Por tanto,

\[
\boxed{
Q_{\text{esfera}}=0
}.
\]

La esfera se encuentra polarizada, pero continúa siendo eléctricamente
neutra.

\end{explicacion}


% ------------------------------------------------------------
% INTERMEDIO 7
% ------------------------------------------------------------

\begin{ejercicio}

Una esfera conductora inicialmente neutra se carga por inducción utilizando
una barra positivamente cargada.

Indique:

\begin{enumerate}
    \item hacia qué lado se desplazan inicialmente los electrones de la
    esfera;
    \item si durante la conexión a tierra entran o salen electrones;
    \item el signo final de la esfera.
\end{enumerate}

\end{ejercicio}

\begin{respuesta}

\begin{enumerate}
    \item Los electrones se desplazan hacia la barra positiva.
    \item Entran electrones desde la Tierra.
    \item La esfera queda cargada negativamente.
\end{enumerate}

\end{respuesta}

\begin{explicacion}

La barra positiva atrae a los electrones móviles de la esfera, por lo que
estos se desplazan hacia la región cercana a la barra.

Cuando la esfera se conecta a tierra mientras la barra permanece próxima,
la atracción producida por la barra favorece la entrada de electrones desde
la Tierra.

Si primero se retira la conexión a tierra, esos electrones adicionales ya no
pueden regresar por ese camino.

Posteriormente se retira la barra y la carga se redistribuye sobre la esfera.

Como el conductor terminó con un exceso de electrones,

\[
\boxed{
Q_{\text{final}}<0
}.
\]

Por tanto, queda negativamente cargado.

\end{explicacion}


% ------------------------------------------------------------
% INTERMEDIO 8
% ------------------------------------------------------------

\begin{ejercicio}

Un estudiante afirma:

\begin{quote}
Si un objeto cargado atrae un trozo de papel, entonces el papel debe tener
carga de signo opuesto.
\end{quote}

¿La afirmación es necesariamente correcta? Justifique.

\end{ejercicio}

\begin{respuesta}

No.

Un cuerpo cargado también puede atraer un objeto eléctricamente neutro debido
a polarización.

\end{respuesta}

\begin{explicacion}

El papel puede ser inicialmente neutro.

Al acercar un cuerpo cargado, las cargas del papel pueden redistribuirse
ligeramente.

Las cargas de signo opuesto al cuerpo externo pueden quedar, en promedio,
más cerca de él, mientras que las cargas del mismo signo quedan más alejadas.

Esta redistribución puede originar una atracción neta aunque

\[
Q_{\text{papel}}=0.
\]

Por tanto, observar atracción no permite concluir necesariamente que los dos
cuerpos poseen cargas netas de signos opuestos.

\end{explicacion}


% ------------------------------------------------------------
% INTERMEDIO 9
% ------------------------------------------------------------

\begin{ejercicio}

Durante un proceso de electrización, un cuerpo \(A\) pierde
\(2.0\times10^{13}\) electrones y estos pasan completamente a un cuerpo
\(B\), inicialmente neutro.

Determine las cargas adquiridas por ambos cuerpos.

\end{ejercicio}

\begin{respuesta}

\[
\boxed{
Q_A\approx+\SI{3.20}{\micro\coulomb}
}
\]

\[
\boxed{
Q_B\approx-\SI{3.20}{\micro\coulomb}
}.
\]

\end{respuesta}

\begin{explicacion}

La magnitud de la carga transferida es

\[
|Q|=Ne.
\]

Sustituyendo,

\[
|Q|
=
\left(2.0\times10^{13}\right)
\left(\SI{1.602176634e-19}{\coulomb}\right).
\]

Entonces,

\[
|Q|
\approx
\SI{3.20e-6}{\coulomb}.
\]

Por tanto,

\[
|Q|
\approx
\SI{3.20}{\micro\coulomb}.
\]

El cuerpo \(A\) perdió electrones, por lo que adquiere carga positiva:

\[
Q_A\approx+\SI{3.20}{\micro\coulomb}.
\]

El cuerpo \(B\) recibió esos mismos electrones, por lo que adquiere carga
negativa:

\[
Q_B\approx-\SI{3.20}{\micro\coulomb}.
\]

La suma de las cargas adquiridas es

\[
Q_A+Q_B=0,
\]

lo cual es consistente con la conservación de la carga.

\end{explicacion}


% ------------------------------------------------------------
% INTERMEDIO 10
% ------------------------------------------------------------

\begin{ejercicio}

Un conductor posee carga neta cero, pero se observa que una región tiene
exceso de electrones y otra región tiene déficit de electrones.

¿Existe una contradicción? Explique.

\end{ejercicio}

\begin{respuesta}

No.

El conductor puede estar polarizado y conservar carga neta igual a cero.

\end{respuesta}

\begin{explicacion}

La carga neta depende de la suma algebraica de todas las cargas del cuerpo.

Una región puede presentar un exceso relativo de electrones y otra un déficit
sin que haya ocurrido transferencia de carga hacia el exterior o desde él.

Por tanto,

\[
Q_{\text{neto}}=0
\]

puede coexistir con una distribución espacial no uniforme de carga.

El cuerpo se encuentra

\[
\boxed{
\text{neutro y polarizado}
}.
\]

No existe contradicción porque neutralidad y polarización describen
propiedades diferentes.

\end{explicacion}


% ============================================================
% NIVEL AVANZADO
% ============================================================

\subsection{Nivel avanzado}

En este nivel se requiere seleccionar principios, analizar sistemas,
identificar información insuficiente y justificar físicamente las
conclusiones.


% ------------------------------------------------------------
% AVANZADO 1
% ------------------------------------------------------------

\begin{ejercicio}

Tres cuerpos aislados \(A\), \(B\) y \(C\) tienen inicialmente una carga
total

\[
Q_{\text{total}}=+\SI{6}{\micro\coulomb}.
\]

Después de varias transferencias internas se mide

\[
Q_A=-\SI{2}{\micro\coulomb},
\qquad
Q_B=+\SI{5}{\micro\coulomb}.
\]

Determine la carga de \(C\) y explique qué principio permite resolver el
problema sin conocer cómo ocurrieron las transferencias.

\end{ejercicio}

\begin{respuesta}

\[
\boxed{
Q_C=+\SI{3}{\micro\coulomb}
}.
\]

El principio utilizado es la conservación de la carga eléctrica.

\end{respuesta}

\begin{explicacion}

El sistema está aislado, por lo que su carga total permanece constante:

\[
Q_{\text{inicial}}
=
Q_{\text{final}}.
\]

La carga total final debe seguir siendo

\[
+\SI{6}{\micro\coulomb}.
\]

Por tanto,

\[
Q_A+Q_B+Q_C
=
+\SI{6}{\micro\coulomb}.
\]

Sustituyendo,

\[
-\SI{2}{\micro\coulomb}
+
\SI{5}{\micro\coulomb}
+
Q_C
=
\SI{6}{\micro\coulomb}.
\]

Entonces,

\[
Q_C
=
+\SI{3}{\micro\coulomb}.
\]

Por tanto,

\[
\boxed{
Q_C=+\SI{3}{\micro\coulomb}
}.
\]

No es necesario conocer las transferencias individuales porque la
conservación fija la carga total del sistema.

\end{explicacion}


% ------------------------------------------------------------
% AVANZADO 2
% ------------------------------------------------------------

\begin{ejercicio}

Se realizan dos observaciones:

\begin{enumerate}
    \item el cuerpo \(A\) repele al cuerpo \(B\);
    \item el cuerpo \(B\) atrae al cuerpo \(C\).
\end{enumerate}

Se sabe que \(B\) está cargado positivamente.

¿Qué puede concluirse con certeza acerca de \(A\) y qué puede concluirse
acerca de \(C\)?

\end{ejercicio}

\begin{respuesta}

\(A\) debe estar cargado positivamente.

No puede afirmarse con certeza que \(C\) sea negativo: puede estar cargado
negativamente o ser neutro y polarizarse.

\end{respuesta}

\begin{explicacion}

Como \(A\) repele a \(B\) y \(B\) es positivo, ambos deben tener cargas del
mismo signo.

Por tanto,

\[
\boxed{
Q_A>0
}.
\]

En cambio, la atracción entre \(B\) y \(C\) admite más de una explicación.

Puede ocurrir porque \(C\) posee carga negativa:

\[
Q_C<0,
\]

pero también puede ocurrir si \(C\) es eléctricamente neutro y se polariza
debido a la presencia de \(B\).

Por ello, la repulsión permite una conclusión más fuerte sobre los signos
que la atracción.

\end{explicacion}


% ------------------------------------------------------------
% AVANZADO 3
% ------------------------------------------------------------

\begin{ejercicio}

Un estudiante propone que una esfera conductora neutra puede cargarse por
inducción siguiendo este procedimiento:

\begin{enumerate}
    \item acercar una barra negativa;
    \item conectar la esfera a tierra;
    \item retirar la barra negativa;
    \item desconectar finalmente la tierra.
\end{enumerate}

¿El procedimiento garantiza que la esfera quede cargada positivamente?
Analice el error.

\end{ejercicio}

\begin{respuesta}

No.

La conexión a tierra debe retirarse antes de retirar la barra inductora.

\end{respuesta}

\begin{explicacion}

Al acercar la barra negativa, los electrones de la esfera son repelidos y el
conductor se polariza.

Cuando se conecta la esfera a tierra, algunos electrones pueden salir hacia
la Tierra debido a la repulsión producida por la barra.

Para que ese déficit de electrones permanezca en la esfera, primero debe
eliminarse el camino conductor hacia la Tierra.

El orden correcto es:

\begin{enumerate}
    \item acercar la barra negativa;
    \item conectar la esfera a tierra;
    \item retirar la conexión a tierra;
    \item retirar finalmente la barra.
\end{enumerate}

Si se retira primero la barra mientras la esfera continúa conectada a tierra,
los electrones pueden volver a fluir entre la Tierra y la esfera.

Por ello, el procedimiento propuesto no garantiza que la esfera conserve una
carga positiva.

\end{explicacion}


% ------------------------------------------------------------
% AVANZADO 4
% ------------------------------------------------------------

\begin{ejercicio}

Dos esferas conductoras aisladas tienen cargas

\[
Q_A=+\SI{10}{\micro\coulomb},
\qquad
Q_B=-\SI{4}{\micro\coulomb}.
\]

Se ponen en contacto.

Un estudiante afirma que después del contacto cada esfera tendrá

\[
+\SI{3}{\micro\coulomb}.
\]

¿La conclusión es siempre válida?

\end{ejercicio}

\begin{respuesta}

No.

Solo puede asegurarse que la carga total del sistema es

\[
\boxed{
Q_{\text{total}}=+\SI{6}{\micro\coulomb}
}.
\]

Para concluir que cada esfera queda con
\(+\SI{3}{\micro\coulomb}\), debe saberse además que son idénticas y están
en condiciones simétricas.

\end{respuesta}

\begin{explicacion}

La conservación de la carga permite calcular la carga total:

\[
Q_{\text{total}}
=
Q_A+Q_B.
\]

Entonces,

\[
Q_{\text{total}}
=
\SI{10}{\micro\coulomb}
-\SI{4}{\micro\coulomb}.
\]

Por tanto,

\[
\boxed{
Q_{\text{total}}
=
+\SI{6}{\micro\coulomb}
}.
\]

Sin embargo, la conservación no determina por sí sola cómo se distribuye esa
carga entre los dos conductores.

La relación

\[
Q_A'=Q_B'
\]

requiere información adicional relacionada con la geometría y la simetría de
los conductores.

Si las esferas fueran idénticas, entonces sí podría escribirse

\[
Q_A'=Q_B'
=
\frac{\SI{6}{\micro\coulomb}}{2}
=
+\SI{3}{\micro\coulomb}.
\]

Sin esa condición, la afirmación no está justificada.

\end{explicacion}


% ------------------------------------------------------------
% AVANZADO 5
% ------------------------------------------------------------

\begin{ejercicio}

Se afirma que durante una electrización por inducción la carga eléctrica no
se conserva porque una esfera inicialmente neutra puede terminar cargada sin
haber tocado al inductor.

Explique por qué el argumento es incorrecto.

\end{ejercicio}

\begin{respuesta}

La carga sí se conserva.

Durante la conexión a tierra, la esfera intercambia electrones con la Tierra,
por lo que la esfera sola no constituye un sistema aislado.

\end{respuesta}

\begin{explicacion}

El principio de conservación de la carga se aplica a un sistema aislado.

Durante la electrización por inducción, la esfera se conecta a tierra.
Mientras existe esa conexión, electrones pueden entrar o salir del conductor.

Por ello,

\[
\text{esfera sola}
\]

no constituye en ese momento un sistema aislado.

La carga que adquiere la esfera procede de una transferencia entre la esfera
y la Tierra.

Para aplicar correctamente la conservación debe considerarse un sistema que
incluya todos los cuerpos entre los cuales puede transferirse carga.

No se creó carga eléctrica; simplemente se produjo una transferencia.

\end{explicacion}


% ------------------------------------------------------------
% AVANZADO 6
% ------------------------------------------------------------

\begin{ejercicio}

Un cuerpo tiene una carga medida de

\[
Q=\SI{7.9e-19}{\coulomb}.
\]

Al dividir entre la carga elemental se obtiene aproximadamente

\[
\frac{Q}{e}\approx4.93.
\]

¿Debe concluirse que la cuantización de la carga es falsa?

\end{ejercicio}

\begin{respuesta}

No.

El resultado es compatible con una carga aproximada de

\[
\boxed{
Q=5e
}
\]

si la medición contiene incertidumbre o ha sido redondeada.

\end{respuesta}

\begin{explicacion}

La cuantización establece

\[
Q=ne,
\]

donde \(n\) debe ser entero.

La razón obtenida experimentalmente es

\[
\frac{Q}{e}\approx4.93.
\]

Este valor está próximo a

\[
n=5.
\]

En una medición real, la carga informada puede estar limitada por la
resolución del instrumento, incertidumbre experimental o redondeo.

Por ello, no debe interpretarse automáticamente \(4.93\) como evidencia de
una carga fundamental fraccionaria.

El valor es compatible con

\[
Q\approx5e.
\]

\end{explicacion}


% ------------------------------------------------------------
% AVANZADO 7
% ------------------------------------------------------------

\begin{ejercicio}

Un estudiante afirma:

\begin{quote}
Un aislante no puede electrizarse porque sus electrones no se mueven.
\end{quote}

Evalúe críticamente la afirmación.

\end{ejercicio}

\begin{respuesta}

La afirmación es incorrecta.

Los aislantes sí pueden electrizarse; la diferencia es que sus cargas poseen
mucha menor movilidad macroscópica y tienden a permanecer localizadas.

\end{respuesta}

\begin{explicacion}

Un aislante también está constituido por partículas cargadas.

La característica que lo distingue de un conductor no es la ausencia de
electrones, sino la menor movilidad de sus portadores de carga.

Un aislante puede ganar o perder electrones durante un proceso de
electrización.

La diferencia es que, una vez depositada, la carga suele permanecer
localizada en la región donde fue transferida.

Por tanto,

\[
\boxed{
\text{aislante}
\not\Rightarrow
\text{imposibilidad de electrización}
}.
\]

\end{explicacion}


% ------------------------------------------------------------
% AVANZADO 8
% ------------------------------------------------------------

\begin{ejercicio}

Una barra negativa se aproxima a una esfera metálica neutra y aislada.

Un estudiante dice:

\begin{quote}
La parte cercana de la esfera se vuelve positiva porque la barra le transmite
carga positiva.
\end{quote}

Explique correctamente el fenómeno.

\end{ejercicio}

\begin{respuesta}

La barra no transmite carga positiva.

Los electrones de la esfera son repelidos hacia la región más alejada,
dejando en la zona cercana un déficit relativo de electrones.

\end{respuesta}

\begin{explicacion}

La esfera comienza siendo eléctricamente neutra.

Al acercar la barra negativa, los electrones móviles de la esfera
experimentan repulsión.

Como consecuencia, se desplazan hacia la región más alejada de la barra.

La zona próxima queda con un déficit relativo de electrones y, por ello,
presenta carga positiva local.

No se produjo transferencia de carga positiva desde la barra.

Además, como ningún electrón salió de la esfera,

\[
Q_{\text{total de la esfera}}=0.
\]

El fenómeno observado es una polarización.

\end{explicacion}


% ------------------------------------------------------------
% AVANZADO 9
% ------------------------------------------------------------

\begin{ejercicio}

Dos cuerpos inicialmente neutros se electrizan por frotamiento. Después del
proceso se mide

\[
Q_A=+\SI{1.5}{\micro\coulomb}
\]

y

\[
Q_B=-\SI{1.2}{\micro\coulomb}.
\]

Si los dos cuerpos fueran realmente el sistema completo y estuvieran
perfectamente aislados durante todo el proceso, ¿son compatibles estas
mediciones con la conservación de la carga?

\end{ejercicio}

\begin{respuesta}

No.

La carga final total sería

\[
\boxed{
Q_{\text{final}}=+\SI{0.3}{\micro\coulomb}
},
\]

aunque inicialmente era cero.

\end{respuesta}

\begin{explicacion}

Como ambos cuerpos eran inicialmente neutros,

\[
Q_{\text{inicial}}=0.
\]

La carga final medida sería

\[
Q_{\text{final}}
=
Q_A+Q_B.
\]

Sustituyendo,

\[
Q_{\text{final}}
=
\SI{1.5}{\micro\coulomb}
-\SI{1.2}{\micro\coulomb}.
\]

Entonces,

\[
Q_{\text{final}}
=
+\SI{0.3}{\micro\coulomb}.
\]

Esto contradice

\[
Q_{\text{inicial}}
=
Q_{\text{final}}
\]

si los dos cuerpos constituyen realmente un sistema perfectamente aislado.

Por tanto, alguna de las hipótesis o mediciones debe revisarse.

Entre las posibilidades están:

\begin{itemize}
    \item intercambio de carga con el entorno;
    \item incertidumbre o error experimental;
    \item existencia de otro cuerpo que participó en la transferencia.
\end{itemize}

La conservación de la carga permite evaluar la coherencia física de las
mediciones.

\end{explicacion}


% ------------------------------------------------------------
% AVANZADO 10
% ------------------------------------------------------------

\begin{ejercicio}

Se dispone de una esfera conductora neutra, una barra negativamente cargada y
una conexión a tierra.

Describa un procedimiento para dejar la esfera con carga positiva
\textbf{sin que la barra toque la esfera}.

Para cada etapa indique qué ocurre con los electrones.

\end{ejercicio}

\begin{respuesta}

El procedimiento es:

\begin{enumerate}
    \item acercar la barra negativa;
    \item conectar la esfera a tierra;
    \item retirar la conexión a tierra manteniendo la barra cerca;
    \item retirar finalmente la barra.
\end{enumerate}

La esfera queda cargada positivamente.

\end{respuesta}

\begin{explicacion}

\textbf{Paso 1: acercar la barra negativa.}

La barra repele a los electrones móviles de la esfera.

Los electrones se desplazan hacia la región más alejada y la esfera se
polariza.

En este momento,

\[
Q_{\text{esfera}}=0.
\]

\textbf{Paso 2: conectar la esfera a tierra.}

La repulsión de la barra negativa favorece que algunos electrones abandonen
la esfera y pasen hacia la Tierra.

La esfera comienza a presentar un déficit de electrones.

\textbf{Paso 3: retirar la conexión a tierra.}

La conexión debe retirarse mientras la barra continúa cerca.

De esta manera se elimina el camino por el cual los electrones podrían volver
desde la Tierra.

\textbf{Paso 4: retirar la barra.}

Al desaparecer la influencia externa, la carga restante se redistribuye
sobre la superficie del conductor.

Como la esfera perdió electrones,

\[
\boxed{
Q_{\text{final}}>0
}.
\]

Por tanto, queda cargada positivamente sin haber tenido contacto directo con
la barra inductora.

\end{explicacion}